\appendix
\section{Installing in a \texttt{conda} environment}
\label{conda_install}

Assuming that \texttt{conda} is installed and activated\ldots

\begin{itemize}

    \item define a new \texttt{conda} environment named {\tt edps}:\newline
    \verb|conda create -n edps python=3.11|

   \item activate the \texttt{edps} environment: \newline
    \verb|conda activate edps|

   \item upgrade {\tt pip}:\newline
   \verb|python -m pip install --upgrade pip|

   \item install EDPS (single line command): \newline
   \verb|python -m pip install --extra-index-url \|\newline
   \verb|   https://ftp.eso.org/pub/dfs/pipelines/repositories/stable/src|\newline
   \verb|   edps edpsgui edpsplot adari_core|

   \item to remove the \texttt{edps} \texttt{conda} environment at anytime:\newline
   \verb|conda env remove --name edps|

\end{itemize}

Use Python 3.11, 3.12, or 3.13 (not 3.14). Then to run \texttt{edps} at any time in the future, activate the \texttt{edps} environment, if it is not already activated, and then simply invoke \texttt{edps}, e.g.\ the first time to create and setup the {\tt application.properties} config file:
\begin{verbatim}
    edps
\end{verbatim}
